% Image credits for the photographs and AI illustrations of Book 5 (University, Year 3).
% Machine-readable license records live in images/book5/CREDITS.md; keep the
% two files in sync. The Book 5 writer owns this file: add one \item per
% photograph (inside the itemize, by chapter) as photographs land.
\chapter*{Image Credits}

The drawings, maps and charts in this book are original; maps and charts
drawn from published data name their source in the caption or in this
list. Photographs are used under their respective free licenses, with
thanks to their authors; full source links and license URLs for every
photograph are listed in \texttt{images/book5/CREDITS.md} in the book's
repository.

\bigskip
Alongside the photographs and the original drawings, some figures are
AI-generated illustrations, produced with OpenAI image generation from
prompts written by the book's editors, and reviewed for scientific
accuracy before inclusion. They show generic scenes only, never a real
place, a real person or a real specimen. The full prompt for every
generated image is recorded in \texttt{images/book5/ai/PROMPTS.md} in the
repository.

\bigskip
\noindent Photographs and data, by chapter:
\begin{itemize}\small
  \item Ch.~1 (seismology and Earth imaging): PREM, A.~M.~Dziewonski and D.~L.~Anderson (1981), as tabulated in ObsPy; ak135, B.~L.~N.~Kennett, E.~R.~Engdahl and R.~Buland (1995).
  \item Ch.~3 (gravity and the figure of the Earth): the twin GRACE Follow-On satellites, artist's concept, NASA/JPL-Caltech, public domain (Wikimedia Commons). Ice-sheet mass balance from the IMBIE Team (2021), UK Polar Data Centre, Open Government Licence v3.0.
  \item Ch.~4 (geomagnetism and palaeomagnetism): International Geomagnetic Reference Field, fourteenth generation, IAGA Division~V Working Group V-MOD, coefficients distributed by NOAA NCEI.
  \item Ch.~5 (thermal models of the lithosphere): age of the ocean floor, M.~Seton and others (2020), EarthByte, CC BY 4.0; ETOPO1 global relief, NOAA NCEI, public domain.
  \item Ch.~6 (mantle convection and plate driving forces): the raised beaches of Ballyhillin, Malin Head, A.~F.~Borchert, CC BY-SA 4.0 (Wikimedia Commons); PREM, A.~M.~Dziewonski and D.~L.~Anderson (1981).
  \item Ch.~7 (deformation mechanisms and shear zones): augen mylonite, Mikenorton, CC BY-SA 4.0 (Wikimedia Commons).
  \item Ch.~9 (extensional and strike-slip tectonics): basin-and-range topography south of Death Valley, Ken Lund, CC BY-SA 2.0 (Wikimedia Commons).
  \item Ch.~14 (stable-isotope geochemistry): LR04 benthic $\delta^{18}$O stack, L.~E.~Lisiecki and M.~E.~Raymo (2005).
  \item Ch.~15 (water--rock interaction): travertine terraces of Palette Spring, Mammoth Hot Springs, Jim Peaco, National Park Service, public domain (Wikimedia Commons). Equilibrium constants from the PHREEQC database phreeqc.dat, D.~L.~Parkhurst and C.~A.~J.~Appelo, U.S. Geological Survey, public domain.
  \item Ch.~16 (sequence stratigraphy): the Book Cliffs above Helper, Utah, David Mackie, public domain (Wikimedia Commons); LR04 benthic $\delta^{18}$O stack, L.~E.~Lisiecki and M.~E.~Raymo (2005).
  \item Ch.~17 (facies models): hummocky cross-stratification in the Arisaig Formation, Michael C.~Rygel, CC BY-SA 3.0 (Wikimedia Commons), cropped.
  \item Ch.~18 (quantitative geomorphology): the Niagara Gorge and the Horseshoe Falls, Ken Lund, CC BY-SA 2.0 (Wikimedia Commons), cropped.
  \item Ch.~19 (geophysical fluid dynamics): zonal-mean winds of the NCEP--NCAR Reanalysis~1, data provided by the NOAA PSL, Boulder, Colorado, USA; World Ocean Atlas 2023, NOAA NCEI.
  \item Ch.~20 (atmospheric dynamics): a low-pressure system off Iceland, NASA/GSFC, MODIS Rapid Response Team, Jacques Descloitres, public domain (Wikimedia Commons); zonal-mean winds of the NCEP--NCAR Reanalysis~1, data provided by the NOAA PSL, Boulder, Colorado, USA.
  \item Ch.~21 (ocean dynamics): an Argo float deployed from a research vessel, NOAA Global Ocean Monitoring and Observing Program, public domain (Wikimedia Commons).
  \item Ch.~22 (climate models): supercomputers of the NASA Center for Climate Simulation, Goddard Space Flight Center, 2010, public domain (Wikimedia Commons); surface air temperature and top-of-atmosphere radiation of the NCEP--NCAR Reanalysis~1, data provided by the NOAA PSL, Boulder, Colorado, USA.
  \item Ch.~24 (glaciology): an outlet glacier in southwest Greenland, NASA / Jim Yungel, Operation IceBridge, 2014, public domain (Wikimedia Commons); Greenland surface elevations from ETOPO1 Ice Surface, NOAA NCEI; Arctic sea-ice extents from the NSIDC Sea Ice Index, version~4 (Fetterer and others, 2025, doi:10.7265/a98x-0f50).
  \item Ch.~25 (remote sensing and GIS): the 1995--2001 interferogram of the Three Sisters volcanic centre, Oregon, U.S. Geological Survey (C.~Wicks and others, 2002), public domain (Wikimedia Commons); laboratory spectra from the ECOSTRESS Spectral Library, version~1.0 (S.~K.~Meerdink, S.~J.~Hook, D.~A.~Roberts and E.~A.~Abbott, 2019; A.~M.~Baldridge and others, 2009), reproduced through the courtesy of the Jet Propulsion Laboratory, California Institute of Technology, Pasadena, California.
  \item Ch.~26 (planetary interiors and surfaces): Barringer Meteor Crater, Arizona, USGS/D.~Roddy, public domain (Wikimedia Commons); Olympus Mons, Viking orbiter mosaic, NASA/JPL/USGS, public domain (Wikimedia Commons). Crater production and chronology functions as given in the Craterstats software of G.~G.~Michael, after G.~Neukum, B.~A.~Ivanov and W.~K.~Hartmann (2001).
  \item Ch.~27 (planetary atmospheres and habitability): a network of dried-up valleys on Mars, Mars Express HRSC, ESA/DLR/FU Berlin, CC BY-SA 3.0 IGO (Wikimedia Commons); Ligeia Mare, Titan, Cassini radar mosaic, NASA/JPL/USGS, public domain (Wikimedia Commons). Habitable-zone fits of R.~K.~Kopparapu and others (2013); vapour pressure of methane from the NIST Chemistry WebBook.
  \item Ch.~28 (exoplanets): HIP~65426~b imaged by the James Webb Space Telescope, NASA/ESA/CSA, A~Carter (UCSC), the ERS~1386 team, and A.~Pagan (STScI), CC BY~4.0 (Wikimedia Commons). Confirmed planets from the NASA Exoplanet Archive, Planetary Systems Composite Parameters table, accessed 23 October 2026: this research has made use of the NASA Exoplanet Archive, which is operated by the California Institute of Technology, under contract with the National Aeronautics and Space Administration under the Exoplanet Exploration Program. Model mass--radius curves of L.~Zeng and others (2016, 2019).
  \item Ch.~29 (the Milky Way and galaxies): the Milky Way panorama of the GigaGalaxy Zoom project, ESO/S.~Brunier, CC BY~4.0 (Wikimedia Commons); the spiral galaxy M101, European Space Agency and NASA (Hubble, heic0602; Hubble data K.~D.~Kuntz, F.~Bresolin, J.~Trauger, J.~Mould and Y.-H.~Chu; image processing Davide De~Martin), CC BY~4.0 (Wikimedia Commons). Circular velocities of the Galaxy from A.-C.~Eilers, D.~W.~Hogg, H.-W.~Rix and M.~K.~Ness (2019); mass model of E.~Pouliasis, P.~Di~Matteo and M.~Haywood (2017); OGLE classical Cepheids of the Large Magellanic Cloud, I.~Soszy\'nski and others (2015), Acta Astronomica 65, 297, distributed by the CDS, Strasbourg (VizieR J/AcA/65/297).
  \item Ch.~30 (where the elements come from): the Crab Nebula, NASA, ESA, J.~Hester and A.~Loll (Arizona State University), public domain (Wikimedia Commons); the kilonova of GW170817 in NGC~4993, NASA and ESA, acknowledgment A.~J.~Levan, N.~R.~Tanvir, A.~Fruchter and O.~Fox (Hubble, heic1717a), CC BY~4.0 (Wikimedia Commons). Solar abundances of M.~Asplund, A.~M.~Amarsi and N.~Grevesse (2021); neutron-capture cross sections from KADoNiS v0.3 (I.~Dillmann, R.~Plag, F.~K\"appeler and T.~Rauscher); half-lives from the IAEA Live Chart of Nuclides (ENSDF); metallicities of the Geneva--Copenhagen survey as re-analysed by L.~Casagrande and others (2011), distributed by the CDS, Strasbourg (VizieR J/A+A/530/A138).
\end{itemize}
